% TOP_PROBLEM: {"id": 2304014, "title": "Research Problems in Function Theory — Problem 4.14", "category_id": 8, "category": {"id": 8, "name": "analysis", "display_name": "Analysis", "description": "Limits, continuity, calculus, and function theory.", "slug": "analysis", "order_index": 8, "created_at": "2026-07-31T15:26:25.671Z"}, "status": "open"}
% FIRST_POSTED: 2026-10-06
% ATTEMPT_STATUS: unresolved
% OPENAI_CLAIM: {"schema_version":1,"source_repo":"https://github.com/openai/math","source_commit":"adc7f1241b42e322a6451854ab7e4b4c146bf78a","release_date":"2026-10-06","match":"partial","resolution":"partial","independently_reviewed":false,"second_pass":"n/a","families":[{"family":"076","title":"Real ultraflat Littlewood polynomials and unbounded binary merit factors","subject":"Real and complex analysis","manuscripts":[{"title":"Nearly minimal maxima and positive minima of Littlewood polynomials","pdf_path":"preprints/Nearly-minimal-maxima-and-positive-minima-of-Littlewood-polynomials-October-5-2026/littlewood-lower-envelope.pdf","readme_path":"preprints/Nearly-minimal-maxima-and-positive-minima-of-Littlewood-polynomials-October-5-2026/README.md","date":"2026-10-05","theorem_ref":"thm:main (sections/introduction.tex:18)","statement_sources":[{"path":"preprints/Nearly-minimal-maxima-and-positive-minima-of-Littlewood-polynomials-October-5-2026/build/sections/introduction.tex","line":18,"label":"thm:main"}]},{"title":"Ultraflat real Littlewood polynomials","pdf_path":"preprints/Ultraflat-real-Littlewood-polynomials-October-5-2026/ultraflat-real-littlewood-polynomials.pdf","readme_path":"preprints/Ultraflat-real-Littlewood-polynomials-October-5-2026/README.md","date":"2026-10-05","theorem_ref":"thm:main (sections/introduction.tex:15)","statement_sources":[{"path":"preprints/Ultraflat-real-Littlewood-polynomials-October-5-2026/build/sections/introduction.tex","line":15,"label":"thm:main"}]}],"lean":{"doc_path":"lean/docs/076.md","comparators":[],"declarations":[],"covers_main_theorem":false},"reasoning_trace":null}]}
\documentclass[11pt]{article}
\usepackage[margin=25mm]{geometry}
\usepackage{fontspec}
\setmainfont{Latin Modern Roman}
\usepackage{amsmath,amssymb,mathtools}
\usepackage{unicode-math}
\setmathfont{Latin Modern Math}
\usepackage{enumitem}
\usepackage{xurl,hyperref,bookmark}
\hypersetup{colorlinks=true,urlcolor=blue}
\setcounter{secnumdepth}{0}
\setlength{\parindent}{0pt}
\setlength{\parskip}{5pt}
\setlength{\emergencystretch}{3em}
\begin{document}
\section*{Research Problems in Function Theory — Problem 4.14}

\subsection{OpenAI's claim}
\textit{Littlewood polynomials with a lower margin}: The paper claims sign polynomials at every sufficiently large length with maximum at most $(1+\eta)\sqrt N$ and minimum at least $\sqrt N/16$, for any fixed positive $\eta$; both inequalities concern the same polynomial.

\textit{Uniformly flat real sign polynomials}: OpenAI claims that real sign polynomials can have normalized modulus uniformly as close to one as desired once the integer length is large enough; choices at different lengths are independent of one another.

\subsection{Mathematical statement of the claimed result}
\paragraph{Mathematical claim 1.} Given $\eta>0$, there is $N_0(\eta)$ such that for every integer $N\ge N_0(\eta)$ one can choose real signs $\epsilon_0,\ldots,\epsilon_{N-1}$ with \[\frac{\sqrt N}{16}\le\left|\sum_{k=0}^{N-1}\epsilon_kz^k\right|\le(1+\eta)\sqrt N\quad(|z|=1).\] The same polynomial has both bounds uniformly on the circle; the lower constant is $1/16$.

\paragraph{Mathematical claim 2.} For every $0<\epsilon<1$, all sufficiently large integers $N$ admit signs $a_k\in\{-1,1\}$ such that \[(1-\epsilon)\sqrt N\le\left|\sum_{k=0}^{N-1}a_kz^k\right|\le(1+\epsilon)\sqrt N\quad\text{for all }|z|=1.\] The quantifier covers every sufficiently large length, and both inequalities refer to one polynomial.

\subsection{Outline of the argument}
\textit{Littlewood polynomials with a lower margin}: The relaxed-coefficient construction for small maxima is augmented on arcs where its leading amplitude is too small. Carefully matched quadratic phases create a positive lower margin without introducing excessive Fourier coefficients outside the permitted degree range. A rounding estimate expressed in terms of the remaining coefficient defect preserves the margin while turning the coefficients into signs.

\textit{Uniformly flat real sign polynomials}: A continuous circle function is built with nearly constant modulus, real coefficients of nearly sign size in the retained frequency range, and a small Fourier tail. Auxiliary-torus waves and phase corrections coordinate the upper and lower modulus bounds. Parseval controls the aggregate coefficient defect, allowing discrepancy rounding to preserve uniform flatness.

\subsection{Scope relative to this problem}
\textit{Littlewood polynomials with a lower margin}: The main theorem supplies a fixed positive lower circle bound and an upper bound at every sufficiently large length (in the two October5 papers). Hayman 4.14 literally asks every n and references omitted constants from 4.13; the small-length and original-source quantifier issue is unresolved, so this broader record is not certified solved.

\textit{Uniformly flat real sign polynomials}: The main theorem supplies a fixed positive lower circle bound and an upper bound at every sufficiently large length (in the two October5 papers). Hayman 4.14 literally asks every n and references omitted constants from 4.13; the small-length and original-source quantifier issue is unresolved, so this broader record is not certified solved.

This page records partial support. Complete resolution of the catalogue question is not established by this record.

\subsection{Formal verification}
Family 076 describes Lean coverage for near-minimal upper maxima and convergence of finite-exponent modulus means; the later uniform lower-envelope conclusions are separate.

\textit{Littlewood polynomials with a lower margin}: No manuscript-specific Comparator is listed in the pinned inventory.

The documented selection does not cover the entire manuscript assertion relevant here.

\textit{Uniformly flat real sign polynomials}: No manuscript-specific Comparator is listed in the pinned inventory.

The documented selection does not cover the entire manuscript assertion relevant here.

This site has not run the Lean development. These links describe OpenAI’s selected formal statements and should be read with their stated scope.

\subsection{Status caveat}
OpenAI reports that an internal model produced the manuscripts. Their arguments have not been independently verified as proofs by this site. Any portion outside the documented formal scope may contain mistakes, and the catalogue label reports the announced mathematical claim.

\subsection{Sources}
\begin{itemize}
\item OpenAI, \emph{Nearly minimal maxima and positive minima of Littlewood polynomials}, 2026-10-05; citation key \texttt{OAI:Nearly-minimal-maxima-and-positive-minima-of-Littlewood-polynomials-October-5-2026}. \href{https://github.com/openai/math/blob/adc7f1241b42e322a6451854ab7e4b4c146bf78a/preprints/Nearly-minimal-maxima-and-positive-minima-of-Littlewood-polynomials-October-5-2026/littlewood-lower-envelope.pdf}{PDF (GitHub)}; \href{https://github.com/openai/math/raw/adc7f1241b42e322a6451854ab7e4b4c146bf78a/preprints/Nearly-minimal-maxima-and-positive-minima-of-Littlewood-polynomials-October-5-2026/littlewood-lower-envelope.pdf}{PDF (direct)}; \href{https://github.com/openai/math/blob/adc7f1241b42e322a6451854ab7e4b4c146bf78a/preprints/Nearly-minimal-maxima-and-positive-minima-of-Littlewood-polynomials-October-5-2026/README.md}{Manuscript page and citation}; \href{https://github.com/openai/math/blob/main/preprints/Nearly-minimal-maxima-and-positive-minima-of-Littlewood-polynomials-October-5-2026/littlewood-lower-envelope.pdf}{Latest version}; Statement source: \href{https://github.com/openai/math/blob/adc7f1241b42e322a6451854ab7e4b4c146bf78a/preprints/Nearly-minimal-maxima-and-positive-minima-of-Littlewood-polynomials-October-5-2026/build/sections/introduction.tex\#L18}{thm:main}.
\item OpenAI, \emph{Ultraflat real Littlewood polynomials}, 2026-10-05; citation key \texttt{OAI:Ultraflat-real-Littlewood-polynomials-October-5-2026}. \href{https://github.com/openai/math/blob/adc7f1241b42e322a6451854ab7e4b4c146bf78a/preprints/Ultraflat-real-Littlewood-polynomials-October-5-2026/ultraflat-real-littlewood-polynomials.pdf}{PDF (GitHub)}; \href{https://github.com/openai/math/raw/adc7f1241b42e322a6451854ab7e4b4c146bf78a/preprints/Ultraflat-real-Littlewood-polynomials-October-5-2026/ultraflat-real-littlewood-polynomials.pdf}{PDF (direct)}; \href{https://github.com/openai/math/blob/adc7f1241b42e322a6451854ab7e4b4c146bf78a/preprints/Ultraflat-real-Littlewood-polynomials-October-5-2026/README.md}{Manuscript page and citation}; \href{https://github.com/openai/math/blob/main/preprints/Ultraflat-real-Littlewood-polynomials-October-5-2026/ultraflat-real-littlewood-polynomials.pdf}{Latest version}; Statement source: \href{https://github.com/openai/math/blob/adc7f1241b42e322a6451854ab7e4b4c146bf78a/preprints/Ultraflat-real-Littlewood-polynomials-October-5-2026/build/sections/introduction.tex\#L15}{thm:main}.
\item Family 076: \href{https://github.com/openai/math/blob/adc7f1241b42e322a6451854ab7e4b4c146bf78a/lean/docs/076.md}{Lean scope}.
\item \href{https://github.com/openai/math/blob/adc7f1241b42e322a6451854ab7e4b4c146bf78a/CONTENTS.md}{OpenAI catalogue entry}.
\end{itemize}
\end{document}
